\documentclass[11pt,a4paper,oneside]{report}
\usepackage{translation}
$if(highlighting-macros)$
$highlighting-macros$
$endif$
\begin{document}
\begin{titlepage}
\color{stcink}
\vspace*{12mm}
{\sffamily\small\color{stcteal}\MakeUppercase{$paper-id$}\par}
\vspace{7mm}
{\sffamily\bfseries\fontsize{25}{31}\selectfont $korean-title$\par}
\vspace{5mm}
{\Large\itshape $original-title$\par}
\vspace{8mm}
{\large $author$\par}
\vfill
\begin{translationnotice}
\textbf{문서 성격}\quad 원문 순서·수치·한정어를 보존한 한국어 구조 보존형 의역본이다.
본문 뒤의 원문 보존 부록은 원 논문의 equation, table, figure, caption을 검증 가능하게 남긴다. 인용과 재현에는 반드시 고정 원문을 사용한다.\\[1mm]
\textbf{고정 버전}\quad $source-version$\\
\textbf{배포 상태}\quad $rights$ — $license-note$
\end{translationnotice}
\vspace{8mm}
{\small\color{stcblue}Sleep-Time Compute Research Program · 2026-08-05\par}
\end{titlepage}
\tableofcontents
\clearpage
$body$
\clearpage
\chapter*{번역·검증 고지}
\addcontentsline{toc}{chapter}{번역·검증 고지}
이 문서는 분석이나 평가를 추가하지 않는 것을 원칙으로 하며, 원문의 주장 강도와 수치를 유지하도록 작성되었다. 한국어 초안 작성과 구조 검수에는 Codex GPT-5.6을 사용했고, 고정 source hash·원문 보존 부록·역대조 표를 통해 검증했다. 불일치가 있을 때에는 원문이 우선한다.
$for(include-after)$
$include-after$
$endfor$
\end{document}
